\documentclass[11pt,a4paper]{article}

\usepackage[T1]{fontenc}
\usepackage[utf8]{inputenc}
\usepackage{textalpha}
\usepackage[spanish]{babel}
\usepackage[a4paper,margin=2.8cm]{geometry}
\usepackage{microtype}
\usepackage{xurl}
\usepackage[hidelinks]{hyperref}

\hypersetup{
  pdftitle={Borges en su biblioteca},
  pdfauthor={Juan José Gómez Cadenas}
}

\setlength{\parindent}{1.25em}
\setlength{\parskip}{0.35em}
\raggedbottom

\title{Borges en su biblioteca}
\author{Juan José Gómez Cadenas}
\date{}

\begin{document}

\maketitle

\begin{quotation}\itshape
La escena se desarrolla durante las altas horas de la noche, en cierta indefinida estancia de una quinta alquilada ---tan indefinida como la estancia excepto por el hecho de su localización, en la calle Gaona, en Ramos Mejía---. Desde el corredor un espejo acecha los movimientos de los dos escritores, repitiendo en la penumbra los objetos que ocupan el cuarto. Sobre la mesa, las sobras de un dulce de leche y una botella mediada de Madeira, que oculta parcialmente el tomo XLVI de la Anglo-American Cyclopaedia, cuyos sucesivos volúmenes se alinean sobre varios anaqueles de la biblioteca, al fondo de la habitación. Uno de estos escritores es el joven Bioy Casares. El otro, Jorge Luis Borges.
\end{quotation}

Pregúntate, sagaz lectora, ojiagudo lector, si el texto que acabas de leer está escrito por un humano. Fíjate en algunos detalles reveladores. El estilo pedante, que copia a Borges hasta en la manera de escupir. Los términos rebuscados, las largas frases. ¿No sospecharías, quizás, la pérfida pluma de una IA vertiendo \emph{slop} digital en las páginas?
Quizás, si has seguido el animado debate entre el sagaz Ángel Fernández y el que suscribe estas líneas, te encuentras tentada, o provocado, a pasar Pangram y ver qué dice. O puede que los argumentos de Ángel y otros escépticos te hayan convencido de que no vale la pena, ya que los detectores de texto generado son chapuceros, poco fiables, fáciles de estafar. Quizás deberíamos simplemente fiarnos de nuestro criterio y lo cierto es que el fragmento que nos ocupa no deja de ser sospechoso.

Pero sigamos. Nuestro texto continúa describiendo la misteriosa Biblioteca de Babel que imaginó el gran Jorge Luis Borges.

\begin{quotation}\itshape
La Biblioteca es una entidad que contiene todo aquello que existe y fuera de la cual nada es imaginable. Las galerías hexagonales rebosan de anaqueles; estos de libros de formato uniforme. De cada galería parte un estrecho corredor que comunica con otra, idéntica a la primera y a todas. Por el centro del corredor pasa la escala espiral que une las diversas plantas de la Biblioteca y que se pierde en el infinito en los dos sentidos. En este corredor, como en la quinta de Ramos Mejía, un espejo duplica fielmente las apariencias. Cada sala hexagonal es, de hecho, la copia exacta de un arquetipo invariable:
\end{quotation}

Sigue una referencia verbatim a Borges:

\begin{quotation}
A cada uno de los muros de cada hexágono corresponden cinco anaqueles; cada anaquel encierra treinta y dos libros de formato uniforme; cada libro es de cuatrocientas diez páginas; cada página, de cuarenta renglones; cada renglón, de unas ochenta letras de color negro.
\end{quotation}

Y otra parrafada, rebosante de locas especulaciones:

\begin{quotation}\itshape
También se nos informa sobre el alfabeto de la Biblioteca (existen veinticinco símbolos gráficos que se enumeran de este modo: veintidós caracteres alfabéticos, la coma, el punto y el espacio) y sobre dos axiomas fundamentales, a saber: en toda la Biblioteca no existen dos libros idénticos y los caracteres gráficos se disponen, en los varios libros, en todas las combinaciones posibles.

No se nos dice cuántos libros contiene la Biblioteca. Es cierto que el narrador da a entender que esta debe ser una cifra gigantesca, pero el relato no explica cuán gigantesca. Acaso, intencionadamente, se están omitiendo hechos que solo unos pocos lectores ---quizás nosotros--- puedan adivinar.

Un elemental cálculo combinatorio nos muestra que el número de volúmenes que esta contiene es $N_V = 25^{N_C}$, donde $N_C$ es el número de caracteres contenidos en un solo libro, es decir, $N_C = 410\times 40\times 80 = 1\,312\,000$, y por lo tanto, $N_V = 25^{1\,312\,000}$. Una simple transformación logarítmica nos permite expresar el resultado en base decimal, $N_V = 10^{\log_{10}(25)\times 1\,312\,000}$. No podemos resistirnos a la tentación de simplificar el número que acabamos de obtener, truncando el exponente al entero más próximo y obteniendo por tanto $N_V = 10^{1\,834\,097}$.
\end{quotation}

¿Tienes ya un veredicto? ¿Te has fijado en la cantidad de dashes, en la forma en que el texto se obsesiona con la numerología, en la insufrible pedantería de la redacción?

Estoy seguro de que ya has decidido. Pero hagamos, por un instante, el esfuerzo de olvidarnos de los detalles de la prosa e intentar entender de qué nos está hablando el autor ---hombre o máquina--- del artículo.

Preguntémonos por el número que acaba de ofrecernos como quien canta un décimo premiado de lotería: $N_V = 10^{1\,834\,097}$. ¿Cuán grande es esa cifra?

 Para hacernos una idea, el número total de estrellas en el Universo es de $10^{22}$. El número total de protones, $10^{80}$; el número de cuantos de luz (fotones), $10^{88}$. Si llenáramos el universo visible de granos de arena, ese número sería $10^{90}$. Todas ellas, cantidades insignificantes comparadas con $N_V$. 
 
 Probemos otra comparación, esta vez con la edad del Universo. Desde el Big Bang hasta el presente han pasado trece mil millones de años, o $1{,}38 \times 10^{10}$ años. Cuando llegue la muerte térmica a la que está abocado, el momento fatídico en el que cada partícula elemental que aún no se haya desintegrado escape del horizonte de sucesos de todas las demás ---es decir, cuando cada electrón, fotón y neutrino del universo se queden solos--- habrán pasado $10^{100}$ años.

Preguntémonos cada cuánto tiempo deberíamos dar cuenta de un libro para poder despacharlos todos antes de ese triste final. El cálculo nos dice que necesitaríamos leer un libro cada $3\cdot 10^{-1\,833\,990}$ segundos, pero eso es imposible, porque el cuanto de tiempo con sentido físico más pequeño disponible en nuestro cosmos cuántico (el tiempo de Planck) es de $5{,}4 \cdot 10^{-44}$~s. Si leyéramos un libro cada intervalo de Planck, aún nos quedarían por leer $10^{1\,834\,097}$ volúmenes. Apenas habríamos gastado la calderilla del número de Babel.

Probemos otra metáfora. Si llenáramos el universo con los libros que en la Biblioteca se almacenan, perfectamente empaquetados, necesitaríamos que este se expandiera mucho más allá de nuestro tiempo actual ($1{,}38 \cdot 10^{10}$ años), hasta llegar a los $10^{16}$ años (diez mil billones de años) para poder encajarlos a todos en el espacio correspondiente. Un espacio, por cierto, tan grande que ni siquiera los fotones de luz, los objetos más veloces de la creación, son capaces de conectar. 

Aunque para entonces no habrá luz en el cosmos. Cuando el Universo llegue a los $10^{16}$ años todas sus estrellas se habrán apagado. Es una imagen un poco espeluznante, ¿no? Un orbe sin luz, sin vida, donde no hay otra cosa que libros.

¿Y qué libros son esos? La respuesta no es menos espeluznante. \emph{Todos los que se puedan escribir}. Así que no es de extrañar que, como nos dice el texto que venimos examinando:

\begin{quotation}\itshape
En una página leemos «Eran los ojos de Alejandro azules como el profundo Egeo». En otra, un nombre de mujer se repite incesantemente. Hay una estantería concreta, en cierta sala en nada diferente de las otras, donde encontramos un libro que contiene cierto hexámetro de \emph{La Eneida}: «Ibant obscuri sola sub nocte per umbram». Más allá ---un inconcebible número de estancias más allá, pero quizás en el mismo libro, apenas unas páginas más tarde--- hallamos otra frase familiar náufraga en el desorden: «La libertad, amigo Sancho, es el más grande de los dones que al hombre dieron los cielos». Perseverando en nuestra búsqueda no está excluido que encontremos mensajes de profetas y de lunáticos, frías enumeraciones mercantiles, ágiles octavas, elaboradas disquisiciones metafísicas, desgarrados sonetos.

Hay en esta ingente Biblioteca libros prácticamente idénticos que solo difieren en un personaje, en cierto detalle trivial, en alguna cifra o dato. También nos es autorizado imaginar que sus libros agotan todas las posibles combinaciones dramáticas. Cristo no muere en la cruz ---Judas ocupa su lugar---. Cartago aplasta a Roma. Héctor derrota a Aquiles frente a las murallas de Troya. Rodrigo Díaz de Vivar no duda en alojarse en el mesón del camino, ignorando las lágrimas de una niña. Dante imagina a Beatriz en el infierno y sitúa a la derecha de Dios a Francesca de Rímini y Paolo Malatesta.
\end{quotation}

¿Pero quizás vale la pena invertir toda la eternidad a cambio de dar con obras maestras nunca antes leídas? ¿No sería ese el sueño de cualquier bibliotecario? ¿No sería ese el sueño de Borges?

Eso, querida lectora, apreciado lector, es lo que el gran escritor se imaginaba. Y cuál no sería su sorpresa cuando, al despertarse en el más allá, fue recibido nada menos que por su amigo Bioy Casares, quien le informó del trabajo que El Gran Bibliotecario le había asignado ---el infeliz Jorge Luis ignora que Bioy no es otra cosa que un \emph{eidolon}, y que el Gran Bibliotecario es un pseudónimo del Gran Satán y la Biblioteca uno de los más profundos círculos del infierno, reservado a escritores, poetas y otros pedantes.

El falso Bioy informa a Borges de que su trabajo es leer todos los libros que pueda desde el instante del Big Bang hasta la muerte térmica del Universo ---\emph{une saison en enfer}--- que abarca $10^{100}$ años. Para acelerar la tarea, además, se le concede el poder de leer un libro cada intervalo de Planck ---no olvidemos que las leyes de la física cuántica rigen incluso en el infierno---. Es decir, Borges puede dar cuenta de un volumen cada $5{,}4\cdot 10^{-44}$~s durante $10^{100}$ años y por tanto puede trasegar nada menos que $10^{151}$ tomos.

Borges se pone a la tarea emocionado, saboreando de antemano el momento en que dará con el soneto más bello del mundo ---uno que supere, por imposible que parezca, al que empieza diciendo: ¿Te acordaste de mí cuando subías / al silencio que sufre la serpiente / prisionera de grillos y de umbrías?---, con la épica que sonrojará a Homero, con la novela que desbancará al \emph{Ulises}, a \emph{Cien años de soledad}, al mismísimo \emph{Quijote}.

Pasan los años. La gran bola de fuego que fue el Universo se enfría, las estrellas se iluminan, el Cosmos se llena de vida. Borges lee y lee, un libro en cada cuanto mínimo de tiempo y no encuentra otra cosa que libros sin sentido. El Universo se expande, empujado por la energía oscura, las estrellas se apagan, los protones se desintegran, los agujeros negros se evaporan, finalmente, fotones, neutrinos y electrones se quedan aislados, cada uno solo en su horizonte de sucesos. Acontece la muerte térmica del Universo y el desesperado escritor no ha encontrado ni un solo libro con sentido.

No puede encontrarlo, de hecho. Un cálculo elemental, basado en la entropía de Shannon ---una letra de un texto con sentido lleva poco más de un bit de información, frente a los 4,6 bits de una letra elegida al azar entre veinticinco---, arroja que el número de libros legibles es descomunal, $10^{513\,000}$, mucho más grande que cualquier cantidad que podamos enumerar en el universo. Pero el número total de libros es monstruosamente mayor, $10^{1\,834\,097}$, y la fracción entre unos y otros, $10^{513\,000 - 1\,834\,097} \simeq 10^{-1\,321\,000} \sim 0$, a todos los efectos, con un exponente negativo más grande que un millón.

Es decir, de nada nos vale ---de nada le vale a Borges--- que la Biblioteca contenga todas las joyas de la literatura universal y todas las variantes de esas joyas. La probabilidad de encontrarse uno de esos libros, incluso de encontrarse un libro con una o dos líneas con sentido, es nula. Borges comprende, finalmente, el terrible castigo que le ha impuesto el Gran Bibliotecario. Pasarse la eternidad leyendo libros sin sentido, secuencias aleatorias de caracteres, basura incomprensible, \emph{slop}.

Permíteme ahora, intrigada lectora, impaciente lector, que desvele el acertijo del texto que venimos leyendo en cursivas. Si te ha parecido texto generado, te equivocabas. El texto es pedante, sí, recargado, lleno de retruécanos, contiene todos los elementos que un crítico mal informado o mal intencionado alegaría que podrían engañar a un detector como Pangram.

Pero fue escrito por dos humanos, concretamente el que suscribe estas líneas y su amigo Marco Guagnelli, hace ya la friolera de 25 años. El texto íntegro lo publicó Jenaro Talens en \emph{Eutopías} en 2021\footnote{\url{https://turia.uv.es/index.php/eutopias/es/article/view/22924}}.

El tono no era casual, ni inocente. Toda la pieza es, a la vez, una parodia y un homenaje a Borges. Y, naturalmente, Pangram asegura que es cien por cien humano. Es muy difícil, por no decir imposible, que en un texto de esa longitud Pangram se equivoque, por muy rebuscado que sea. Y creo que vale la pena que explique por qué.

Primero, un par de ideas muy sencillas. Los humanos escribimos de manera bastante impredecible para una máquina (o al menos lo hemos venido haciendo hasta ahora). Un LLM está entrenado para escoger siempre la palabra más probable (o una de las más probables) que continúa con el texto que tiene en contexto. Así que si un LLM lee el texto generado por otro LLM, se da cuenta enseguida de que las palabras que este genera para cada longitud de texto dado lo sorprenden poco. Él haría lo mismo. Si lee uno de mis textos, se sorprende cada poco. Como cualquier escritor con estilo propio, salgo a menudo por peteneras y me voy por los cerros de Úbeda. Así que la máquina se sorprende muy a menudo y esa sorpresa se puede medir y recibe, de hecho, un nombre muy sugestivo. \emph{Perplejidad}. Cuanto más perpleja se queda la máquina, más humano el texto.

El problema es que la perplejidad, al final, cuantifica una cierta probabilidad basándose en un único criterio ---el de la palabra más probable que sigue al texto--- y eso hace posible que se pueda entrenar a un LLM para que cambie sus patrones, eligiendo con mucha más frecuencia palabras improbables de lo que lo haría en condiciones normales para engañar a un detector basado en medir perplejidad. Y viceversa, textos rebuscados como el que he mostrado aquí se parecen demasiado a fragmentos «borgianos» y podrían llevar al detector a no sorprenderse.

Esa es la razón por la que la nueva generación de detectores de IA como Pangram no se basan en una variable, sino que entrenan un LLM ---una bestia de millones de parámetros, similar a los modelos que conocemos, como Claude o ChatGPT--- para que distinga entre la forma de escribir de humanos y máquinas. Como buen LLM, el truco está en poder entrenarlo con suficiente texto 100~\% humano y con suficiente texto 100~\% IA. Y eso no es problema. Disponemos de todas las obras escritas antes de 2022, desde la \emph{Odisea} hasta el \emph{Ulises}, pasando por \emph{El Quijote}, Shakespeare, Goethe y Víctor Hugo, para entrenar al detector ---y de infinitos informes burocráticos, noticias de periódicos, artículos científicos, etcétera---. También podemos usar todo el texto IA que haga falta. Los ingenieros de Pangram pueden generar tanto texto sintético como quieran. Es más, pueden hacerlo usando Claude, ChatGPT y cualquier otro modelo. Para cada texto humano pueden generar (y así lo hacen) miles de variantes de ese mismo texto, con diferentes prompts y diferentes modelos y en cada caso entrenar al LLM diciéndole qué es humano y qué no. El resultado es que los pesos del LLM capturan, por supuesto, las diferencias en perplejidad y ritmo, pero también miles, o decenas de miles de otros aspectos que no pueden describirse fácilmente. En eso no son diferentes de cualquier otro LLM o de otros sistemas IA, como un detector de visión (en este caso una red convolucional) que sabe distinguir, digamos, entre perros y gatos basándose en millones de detalles capturados en sus pesos (en lugar de usar características, como las orejas o los bigotes, como hacían los antiguos sistemas) o sistemas «sobrehumanos» como AlphaZero, que «capturan» la capacidad de jugar al ajedrez con un Elo por encima de 3.500 sin que nadie les tenga que dar una sola lección de táctica o estrategia (como a los antiguos motores de ajedrez).

Por más que se empeñen los prestidigitadores, sus trucos para hacerle saltar los fusibles al detector no prueban nada. Se basan en manipulaciones arteras ---como introducir caracteres extraños--- \emph{en textos cortos}. Todos los detectores de IA avisan de su falta de fiabilidad en textos inferiores a unas cien palabras. Pero si un texto de 10.000 palabras ---no digamos ya una novela larga--- da rojo en Pangram, es que lo ha generado una IA. Puede darse el caso, con suficiente paciencia y en textos lo bastante formales, de que alguien parta de un texto IA y a base de ediciones acabe por conseguir un texto «humano» ---hubiera sido más fácil para el tramposo escribirlo directamente---. También es posible partir de una base IA y que el detector siga cantando aún después de muchas ediciones. En todos los casos, se trata de ejercicios inútiles.

Cualquier lector puede comprobarlo por sí mismo, mandando al detector sus propios textos y certificando que siempre salen humanos. También puede hacer lo contrario y ver cómo un texto IA da rojo chillón. Ninguno de los escépticos que conozco me ha mostrado todavía, ni creo que lo haga, un texto publicado antes de 2022 que Pangram marque como IA. Pangram funciona bien por la misma razón que funcionan bien Claude y ChatGPT. Porque es un potente LLM, entrenado con una ingente cantidad de datos.

Pero, argumentará algún escéptico, el estilo de escritura cambia y quizás dentro de unos años la escritura humana se parecerá tanto a la de los LLM, que Pangram y el resto de detectores darán falsos positivos.

Esa, quizás, es la razón más importante para defender el uso de detectores de IA. Porque lo malo no es que un autor o dos usen la IA para escribir sus novelas, o algún bloguero haga lo propio en Substack, o haya quien tire de máquina en Jot Down. Lo malo es que llegue el día en que todo el mundo lo haga, que llegue el día en que todo el mundo escriba igual, que llegue el día en que todo el mundo piense igual. 

\emph{Ese día, como Borges, estaremos atrapados en una inmensa Biblioteca llena de slop}.

Cuando todos usemos la IA para escribir nuestros textos, habremos delegado escribir a los LLM. Pero escribir es pensar. El día que dejemos de escribir, el día en que nos dé pereza expresarnos, el día en que le pidamos a la máquina que redacte nuestros artículos científicos, nuestro blog y nuestras cartas de amor, nos habremos condenado al infierno.

\end{document}
